\needspace{0.25\textheight}
\section{Численная верификация}\label{sec:numerics}

\subsection{Протокол}

Все проверки выполняются интегратором Рунге--Кутты четвёртого порядка
с постоянным шагом $h$ в двойной точности:
\begin{gather*}
  y_{k+1} = y_k + \frac{h}{6}\bigl(k_1 + 2k_2 + 2k_3 + k_4\bigr),\\
  \begin{aligned}
    k_1 &= f(t_k, y_k), &\quad k_2 &= f\bigl(t_k + \tfrac h2,\; y_k + \tfrac h2 k_1\bigr),\\
    k_3 &= f\bigl(t_k + \tfrac h2,\; y_k + \tfrac h2 k_2\bigr), &\quad k_4 &= f(t_k + h,\; y_k + h k_3),
  \end{aligned}
\end{gather*}
порядок сходимости контролируется удвоением шага (Ричардсон):
$e(h)/e(h/2) = 2^{p}$, $p \approx 4$. Исполнение организовано через
суверенный движок \texttt{poler-engine~0.61.0} (режим
\texttt{--exec}): скрипт верификации запускается движком как единая
команда, а скалярные кросс-проверки ниже в подразделе~\ref{sec:calc}
выполнены его встроенным калькулятором (решатель полиномов
Дюрана--Кернера). Среда: numpy~2.1.3, детерминированный датчик
(seed $= 20261002$), полное время прогона $47$~с.

Символьные параметры документа конкретизируются так:
$\bbI = \diag(0{,}90;\,0{,}62;\,0{,}37)$ (общее положение),
$\Om_0 = (2{,}2;\,-1{,}1;\,0{,}7)$; для диссипативных тестов
$\nu = 0{,}10$, $\bbL = \diag(1;\,1{,}7;\,2{,}3)$; осесимметрия
$I_1 = I_2 = 1$, $I_3 = 0{,}5$, $\bbL = \diag(1;1;2)$, $\nu = 0{,}15$;
контрпример замечания~\ref{rem:commutator}~---
$\bbI^{-1} = \diag(10;1;1)$,
$\bbL = \bigl(\begin{psmallmatrix}1&1&0\\1&1{,}5&0\\0&0&1\end{psmallmatrix}\bigr)$;
Бюргерс: $\nu = 0{,}05$ и $0{,}10$, $u_0 = \sin x$,
$\phi_0 = \exp\bigl(-(1-\cos x)/(2\nu)\bigr)$, $N = 2048$ гармоник,
$\phi$ эволюционирует спектрально точно.

\subsection{Результаты}

\begin{longtable}{@{}l >{\raggedright\arraybackslash}p{6.8cm} l l@{}}
\caption{Верификация модели: 23 проверки, все пройдены. Допуски указаны
в описаниях; шаги и горизонты~--- в протоколе и описаниях.}\label{tab:verify}\\
\toprule
ИД & Проверка (допуск) & Измерено & Статус \\
\midrule
\endfirsthead
\caption[]{(продолжение таблицы~\ref{tab:verify})}\\
\toprule
ИД & Проверка (допуск) & Измерено & Статус \\
\midrule
\endhead
\bottomrule
\endfoot
C1 & Дрейф Казимира $C=\|\Om\|^2$, $\nu=0$, $h=10^{-2}$, $T=200$ \ ($\le 2\cdot10^{-8}$, $p\in[3{,}5;5{,}5]$) & $5{,}73\cdot10^{-9}$, $p=4{,}96$ & PASS \\
C2 & Дрейф гамильтониана $\mathcal H$, $\nu=0$ \ ($\le 2\cdot10^{-8}$) & $7{,}50\cdot10^{-9}$ & PASS \\
P1 & Замкнутость орбиты: $d(T_p)/\|\Om_0\|$ \ ($\le 10^{-5}$) & $4{,}98\cdot10^{-15}$, $T_p=3{,}2603$ & PASS \\
D1 & Невязка $\dot C = -2\nu\langle\Om,\bbL\Om\rangle$ \ ($\le 5\cdot10^{-4}$) & $7{,}75\cdot10^{-5}$ & PASS \\
D2 & Монотонность $C(t)$ при $\nu>0$ \ ($\le 0$) & $0{,}00$ & PASS \\
D3 & $C(t) \le C_0e^{-2\nu\lambda_{\min}t}$ \ ($\le 1+10^{-7}$) & $1{,}000000$ & PASS \\
D4 & Монотонность $\mathcal H$ при \textup{(H2)} \ ($\le 0$) & $0{,}00$ & PASS \\
E1a & Контрпример \ref{rem:commutator}: $\lambda_{\min}(\bbL)$ \ ($>0$) & $0{,}2192$ & PASS \\
E1b & $\dot{\mathcal H}(0) = -\nu\langle\Om,\bbI^{-1}\bbL\Om\rangle$ \ ($=+0{,}15$) & $+0{,}1500$ & PASS \\
E1c & Рост $\mathcal H$ при строго убывающем $C$ & $\mathcal H_{\max}/\mathcal H_0 = 1{,}0102$ & PASS \\
B1 & $\sup\|\Om(t)\|/\|\Om_0\|$, $\nu=0$, $T=1000$ \ ($\le 1+10^{-7}$) & $1{,}0000000002$ & PASS \\
B2 & $\|\Om(t)\| \le \|\Om_0\|e^{-\nu\lambda_{\min}t}$, $T=60$ \ ($\le 1+10^{-9}$) & $1{,}0000000000$ & PASS \\
B2b & $\|\Om(T)\|/\|\Om_0\|$, $T=60$ \ ($\le e^{-6}\approx 2{,}48\cdot10^{-3}$) & $2{,}15\cdot10^{-3}$ & PASS \\
A1 & Явное решение \eqref{eq:exact-nu0}, $h=0{,}25$ \ ($p=4\pm0{,}3$) & $5{,}64\cdot10^{-3}$, $p=3{,}97$ & PASS \\
A2 & Явное решение \eqref{eq:exact-nup}, $\nu=0{,}15$ \ ($p=4\pm0{,}3$) & $2{,}40\cdot10^{-18}$, $p=3{,}99$ & PASS \\
M1 & Точный проектор: $\|\bbP^2-\bbP\|$, $\|\mathcal Q(\bbP)-\bbP\|$ \ ($\le 10^{-13}$) & $2{,}7\cdot10^{-16}$ / $8{,}9\cdot10^{-16}$ & PASS \\
M2 & Острота \eqref{eq:q-bound}: $\max\, \delta_Q/(4\delta_P^2)$, $10^4$ возмущённых $\bbP$ \ ($\le 1$) & $0{,}999997$ & PASS \\
H1a & Невязка Бюргерса, $\nu=0{,}1$: абс./отн. \ (отн. $\le 10^{-3}$) & $8{,}92\cdot10^{-6}$ / $6{,}55\cdot10^{-6}$ & PASS \\
H1b & $\max|u_x|$ при $t=0;\,0{,}5;\,1;\,2;\,5$ \ ($<10$; убывание после $t\approx1$) & $1{,}00/1{,}83/4{,}43/7{,}96/2{,}48$ & PASS \\
S1a & Интегрируемый штрейн $s_0e^{-t}$: $\max\|\Om\|/\mathrm{bound}$ \ ($\le 1+10^{-9}$) & $1{,}0000000000$ & PASS \\
S1b & Изотропный $S=s_0\mathrm{Id}$: $\mathrm{ratio}(T)$, рост $2{,}56 \to 139{,}63$ = границе \ ($=1$) & $0{,}99999992$ & PASS \\
L1a & $2D$-сечение: $\max|\Omega_3(t)|$ \ ($\le 10^{-12}$) & $0{,}00$ & PASS \\
L1b & Невязка $\dot C_2 = -2\nu L_{11}(\Omega_1^2+\Omega_2^2)$ \ ($\le 5\cdot10^{-4}$) & $1{,}00\cdot10^{-4}$ & PASS \\
\end{longtable}

\subsection{Скалярные кросс-проверки калькулятором движка}\label{sec:calc}

Два ключевых скалярных факта раздела~\ref{sec:cascade} проверены
встроенным калькулятором \texttt{poler-engine} независимо от
Python-конвейера:
\[
  \texttt{calc\; 3*0.5\^{}2 - 2*0.5\^{}3} \;\longrightarrow\; 0{,}5
  \qquad\text{(неподвижная точка } q(\tfrac12)=\tfrac12\text{),}
\]
\[
  \texttt{calc\; 0.25\^{}2*(3-1)*(2*0.5+1)} \;\longrightarrow\; 0{,}25
  \;=\; 4\cdot(0{,}25)^2
  \qquad\text{(равенство в \eqref{eq:q-bound} при } \lambda=\tfrac12\text{).}
\]
Вторая проверка~--- числовое подтверждение оптимальности константы $4$
(теорема~\ref{th:mcweeny}): левая и правая части оценки совпадают.

\subsection{Комментарии к таблице}

\emph{Инварианты (C1--C2, P1).} Дрейф Казимира и гамильтониана при
$\nu = 0$ имеет порядок $O(h^4)$: измеренный порядок $p = 4{,}96$
подтверждает асимптотику RK4, а замкнутость орбиты
$d(T_p)/\|\Om_0\| = 5\cdot10^{-15}$~--- машинная точность~---
реализует утверждение предложения~\ref{prop:periodic}: численная
траектория возвращается в начальную точку с погрешностью на уровне
двойной точности.

\emph{Диссипация и тонкость (D1--E1).} Тождество
\eqref{eq:casimir} выполняется вдоль траектории с невязкой
$7{,}8\cdot10^{-5}$ (предел конечно-разностной проверки), монотонность
$C$ и $\mathcal H$~--- с нулевым превышением. Контрпример
замечания~\ref{rem:commutator} воспроизведён точно: аналитическое
$\dot{\mathcal H}(0) = +0{,}15$ совпадает с расчётным до $10^{-12}$,
и гамильтониан вырастает на $1{,}02\%$ при строго убывающем Казимире~---
коммутационная гипотеза \textup{(H2)} действительно необходима.

\emph{МакВини и Хопф--Коул (M1--H1).} Точный ранг-1 проектор модели
идемпотентен до $2{,}7\cdot10^{-16}$, а максимальное отношение
$\delta_Q/(4\delta_P^2) = 0{,}999997$ по $10^4$ случайных возмущённых
проекторов (спектр в максимуме $\{0;\,0{,}5017;\,1\}$) подтверждает
и справедливость, и остроту оценки \eqref{eq:q-bound}. Невязка
тождества Коула--Хопфа при $\nu = 0{,}1$ составляет $6{,}6\cdot10^{-6}$
относительно величин слагаемых (усиление округления через $1/\phi$ с
динамическим диапазоном $e^{1/(2\nu)} = e^5$); поведение
$\max|u_x|$ демонстрирует переформирование фронта до $t\approx1$ с
последующим вязким сглаживанием~--- глобальную регулярность
следствия~\ref{cor:burgers-global}.

\emph{Границы метода (S1, B1--B2).} Оценка Грёнуолла
теоремы~\ref{th:strain} выполняется с нулевым нарушением при
интегрируемом штрейне и достигается \emph{точно} при изотропном
$S = s_0\mathrm{Id}$: рост $\|\Om\|$ от $2{,}56$ до $139{,}63$
совпадает с границей во всех знаках~--- константа в \eqref{eq:strain-bound}
неулучшаема, и без условия интегрируемости штрейна экспоненциальный
рост реализуется. Тем самым численно воспроизведена точная граница
метода, установленная в \S\ref{sec:critique}.

Полный код верификации~--- самодостаточный скрипт
\texttt{verify\_\allowbreak so3.py} (Python/numpy, без внешних зависимостей),
поставляемый вместе с документом; трасса прогона~---
\texttt{verification\_\allowbreak output.txt}. Скрипт детерминирован
(фиксированный seed) и воспроизводит таблицу~\ref{tab:verify} одной
командой.
